\chapter{Как мир попадает внутрь}
% полная версия: chapters/11_sensors.tex

Любой орган чувств --- это преобразователь, как микрофон или фотоэлемент. Свет, звук или молекула действуют на особый белок в клетке-датчике, белок открывает «кран», течёт ток, и в конце цепочки получаются щелчки, которые уходят в сеть. Почти все датчики выбрасывают глутамат, так что входы машины подключаются прямо к её телефонной сети.

\section*{На пределе физики}

Датчики мозга работают у самого предела возможного. Датчик света в темноте отвечает на одну-единственную частицу света. Датчик звука замечает смещение меньше нанометра --- меньше размера нескольких атомов. В полутьме точность ограничивает уже не глаз, а сам свет: частицы приходят случайно, и чтобы различить оттенок, их нужно накопить побольше. Поэтому в сумерках мы видим медленнее --- глаз дольше копит свет, как фотоаппарат с длинной выдержкой.

\section*{Автоматическая экспозиция}

Освещённость от звёздной ночи до солнечного снега меняется в миллиарды раз, громкость звука --- ещё сильнее. А частота щелчков меняется от нуля до нескольких сотен в секунду. Втиснуть одно в другое напрямую невозможно. Решение то же, что у фотоаппарата: автоматическая экспозиция. Датчик всё время подстраивает чувствительность под средний уровень вокруг и сообщает не яркость, а яркость \emph{по сравнению с фоном}. Постоянное со временем перестаёт вызывать ответ, а каждое изменение --- вызывает. Входы машины с самого начала говорят о переменах.

\pencilfig{125}{%
% light rises in two tenfold steps, then drops a hundredfold; sensor rate = baseline + gain * (log light - adapted level),
% the adapted level follows log light with a 0.7 s time constant; rate clipped at zero
\draw[penlite=pcAmber] (0,1.75) -- (1.4,1.75) -- (1.4,2.1) -- (4.2,2.1) -- (4.2,2.45) -- (6.3,2.45) -- (6.3,1.75) -- (8.4,1.75);
\node[lbl, text=pcAmber!70!black, anchor=east] at (-0.05,2.0) {свет};
\node[lbl, anchor=south] at (2.8,2.12) {ярче в 10 раз}; \node[lbl, anchor=south] at (5.25,2.47) {ещё в 10};
\node[lbl, anchor=south] at (7.35,1.77) {снова темно};
\draw[pen] (0,0) -- (8.4,0);
\draw[pcGray, densely dashed, line width=0.5pt] (0,0.32) -- (8.4,0.32);
\draw[penlite=pcBlue] plot coordinates {(0.00,0.32) (1.26,0.32) (1.40,1.02) (1.54,0.85) (1.68,0.72) (1.82,0.62) (1.96,0.54) (2.10,0.49) (2.24,0.45) (2.38,0.41) (2.52,0.39) (2.66,0.37) (2.80,0.36) (3.08,0.34) (3.50,0.33) (4.06,0.32) (4.20,1.03) (4.34,0.85) (4.48,0.72) (4.62,0.62) (4.76,0.54) (4.90,0.49) (5.04,0.45) (5.18,0.41) (5.32,0.39) (5.46,0.37) (5.60,0.36) (5.88,0.34) (6.16,0.33) (6.30,0.00) (7.00,0.00) (7.14,0.07) (7.28,0.13) (7.42,0.18) (7.56,0.22) (7.70,0.24) (7.84,0.26) (7.98,0.28) (8.12,0.29) (8.40,0.30)};
\node[lbl, text=pcBlue, anchor=east] at (-0.05,0.32) {ответ};
\node[lbl, anchor=north] at (6.65,-0.02) {молчит};
}{Свет становится ярче ступенями, каждая в десять раз. Датчик отвечает на каждую перемену всплеском и быстро возвращается к прежнему уровню. Когда снова темнеет, он на время замолкает.}

Инженеры повторили этот приём в событийных камерах. Такая камера не снимает кадры по часам: каждый её пиксель сам сообщает «стало ярче» или «стало темнее», когда что-то изменилось. Это та самая пара проводов «да» и «нет» из второй главы, только в кремнии.

\section*{Глаз сжимает картинку}

В глазу около ста миллионов датчиков света, а проводов, которые везут изображение в мозг, --- около миллиона. Прежде чем картинка покинет глаз, её сжимают примерно в сто раз. Ровные участки почти ничего не стоят, передаются края и перемены. По оценке, сделанной на глазах животных и пересчитанной на человека, глаз передаёт в мозг порядка десяти миллионов бит в секунду --- как старый сетевой кабель.

\section*{Ухо --- это пианино}

Ухо раскладывает звук по частотам, как инженер поставил бы набор фильтров. Упругая перегородка в ухе в разных местах откликается на разные частоты, и датчики на каждом участке слушают свою «клавишу». Номер сработавшего датчика --- это высота тона.

\pencilfig{11}{\pencilart{11}}{Ухо устроено как пианино наоборот: каждая клавиша слушает свою ноту.}

Более того, ухо не пассивно. Часть его клеток сама раскачивает перегородку в такт звуку и усиливает тихие звуки в тысячи раз, а громкие --- почти не усиливает. Это усилитель у самой границы самовозбуждения: там он особенно чувствителен. А на низких частотах нейроны ещё и щёлкают в такт волне, сохраняя точное время звука, --- по нему мозг и находит направление на источник. Чем выше частота, тем хуже нейроны за ней успевают, и остаётся только код местом.

\section*{Остальные чувства}

Почти в каждом чувстве узнаётся приём из предыдущих глав. Температуру передают два провода, отдельные для тепла и для холода. Осязание пользуется датчиками, которые отвечают на прикосновение и сразу замолкают, и датчиками, которые держат ответ, пока давление есть. А обоняние устроено как штрихкод: типов обонятельных датчиков у человека около четырёхсот, и запах --- это набор сработавших типов. Число таких наборов астрономическое.

\fullversion{11}
